%% ITEM 0 TYPE Section
\section{Introduction}%% ITEM 1 TYPE Frame
\begin{frame}
\frametitle{LLM Post-Training Paradigms}
  \begin{itemize}
    \item \textbf{Supervised Fine-Tuning (SFT)}
      \begin{itemize}
        \item Trains on curated input-output pairs
        \item Aligns model behavior with human demonstrations
      \end{itemize}
    \item \textbf{Reinforcement Learning (RL)}
      \item Optimizes policies using feedback signals (e.g., rewards)
      \item Requires environment interaction for policy improvement
    \item \textbf{Key Limitations}
      \begin{itemize}
        \item SFT lacks adaptability to dynamic environments
        \item RL demands extensive interaction/data collection
      \end{itemize}
  \end{itemize}
\end{frame}%% ITEM 2 TYPE Section
\section{Unified Framework}%% ITEM 3 TYPE Frame
\begin{frame}
\frametitle{Unified Policy Gradient Estimator: Overview}
\begin{itemize}
  \item Unified framework for gradient calculations in post-training methods
  \item Combines \textbf{Supervised Fine-Tuning (SFT)} and \textbf{Reinforcement Learning (RL)} approaches
  \item Provides closed-form policy gradients for diverse algorithms
  \item Decomposes methods into interpretable components
  \item Enables systematic comparison via Table \ref{tab:grad_compare} (see next slide)
\end{itemize}
\end{frame}%% ITEM 4 TYPE Frame
\begin{frame}
\frametitle{Key Components of the Estimator}
\begin{itemize}
  \item \textbf{Policy Network}: Core decision-making model parameterized by $\theta$
  \item \textbf{Value Function}: Estimates expected cumulative reward $V(s)$
  \item \textbf{Advantage Estimation}: Measures state-action value relative to baseline $A(s,a)$
  \item \textbf{Baseline Subtraction}: Reduces variance while maintaining unbiasedness
  \item \textbf{Entropy Regularization}: Encourages exploration through policy entropy $H(\pi_\theta)$
\end{itemize}
\end{frame}%% ITEM 5 TYPE Frame
\begin{frame}
\frametitle{Unification in Practice}
\begin{itemize}
  \item Bridges SFT and RL through shared gradient formulation
  \item Flexible architecture adapts to different reward structures
  \item Reduces variance compared to isolated methods
  \item Enables hybrid training approaches (SFT + RL)
  \item Provides theoretical guarantees for convergence
\end{itemize}
\end{frame}%% ITEM 6 TYPE Frame
\begin{frame}
\frametitle{Unification in Practice}
\begin{itemize}
  \item Bridges SFT and RL through shared gradient formulation
  \item Flexible architecture adapts to different reward structures
  \item Reduces variance compared to isolated methods
\end{itemize}
\end{frame}%% ITEM 7 TYPE Frame
\begin{frame}
\frametitle{Unified Policy Gradient Estimator: Overview}
\begin{itemize}
  \item Unified framework for gradient calculations in post-training methods
  \item Combines \textbf{SFT} and \textbf{RL} approaches
  \item Provides closed-form policy gradients
  \item Decomposes methods into interpretable components
\end{itemize}
\end{frame}%% ITEM 8 TYPE Frame
\begin{frame}
\frametitle{Unified Policy Gradient Estimator: Summary}
\begin{itemize}
  \item General framework for policy gradient calculations
  \item Combines SFT and RL through unified gradient formulation
  \item Reduces variance and enables hybrid training methods
  \item Provides interpretable decomposition of algorithms
  \item Theoretical guarantees for convergence and stability
\end{itemize}
\end{frame}%% ITEM 9 TYPE Frame
\begin{frame}
\frametitle{Unified Policy Gradient Estimator: Summary}
\begin{itemize}
  \item General framework for policy gradient calculations
  \item Combines SFT and RL through unified gradient formulation
  \item Reduces variance and enables hybrid training methods
  \item Provides interpretable decomposition of algorithms
  \item Theoretical guarantees for convergence and stability
\end{itemize}
\end{frame}%% ITEM 10 TYPE Frame
\begin{frame}
\frametitle{Unified Policy Gradient Estimator}
\begin{itemize}
  \item Unified framework combining SFT and RL
  \item Provides closed-form policy gradients
  \item Decomposes into interpretable components
  \item Reduces variance through baseline subtraction
  \item Theoretically guaranteed convergence
\end{itemize}
\end{frame}%% ITEM 11 TYPE Frame
\begin{frame}
\frametitle{Unified Policy Gradient Estimator: Core Concepts}
\begin{itemize}
  \item Single framework for SFT and RL gradient calculations
  \item Closed-form expressions for policy gradients
  \item Component-based decomposition (policy, value, advantage, baseline, entropy)
  \item Variance reduction through baseline subtraction
  \item Theoretical convergence guarantees
\end{itemize}
\end{frame}%% ITEM 12 TYPE Frame
\begin{frame}
\frametitle{Unified Policy Gradient Estimator: Summary}
\begin{itemize}
  \item General framework for policy gradient calculations
  \item Combines SFT and RL through unified gradient formulation
  \item Reduces variance and enables hybrid training methods
  \item Provides interpretable decomposition of algorithms
  \item Theoretical guarantees for convergence and stability
\end{itemize}
\end{frame}%% ITEM 13 TYPE Frame
\begin{frame}
\frametitle{Unified Policy Gradient Estimator}
\begin{itemize}
  \item Unified framework combining SFT and RL
  \item Closed-form policy gradient expressions
  \item Component-based decomposition:
    \begin{itemize}
      \item Policy network $\pi_\theta$
      \item Value function $V(s)$
      \item Advantage estimation $A(s,a)$
    \end{itemize}
  \item Variance reduction through baseline subtraction
  \item Theoretical convergence guarantees
\end{itemize}
\end{frame}%% ITEM 14 TYPE Frame
\begin{frame}
\frametitle{Unified Policy Gradient Estimator}
\begin{itemize}
  \item Unified framework combining SFT and RL
  \item Closed-form policy gradient expressions
  \item Key components: policy network, value function, advantage estimation
  \item Variance reduction through baseline subtraction
  \item Theoretical convergence guarantees
\end{itemize}
\end{frame}